\section{Uniform compact-unitary dimensions and an independent trace computation}
\label{sec:kt-unitary}

All statements here use the standard positive Hermitian form on
$E^3$, where $E=\mathbf Q(\omega)$ and $\omega^2+\omega+1=0$,
and the full stabilizer $K$ of $O_E^3$ at the finite places.
Write $V_d$ for the compact $U(3)$ representation of highest weight
$(d,0,-d)$ and $M_d=M(U(3),K,V_d)$.
The lattice class-number-one statement and
$\Gamma=U(3)(\mathbf Q)\cap K=\mu_6^3\rtimes S_3$ are carried
from K17-U3-HECKE, proof paragraphs ``The lattice class number'' and
``The invariant harmonic polynomial'', equations (UH.2)--(UH.4).
In particular, $M_d=V_d^\Gamma$ and $|\Gamma|=1296$.

\subsection{The necessary weight correction}
\begin{proposition}[KT-U3-WEIGHT]
The coefficient space used by K17 for the Selmer exponent $e$ is
$M_{e-1}$, whose highest weight is $(e-1,0,1-e)$.
Its Harish--Chandra parameter is $(e,0,-e)$.
Consequently the number $157899$ belongs to coefficient weight
$(5,0,-5)$ and parameter $(6,0,-6)$.
The literal coefficient space of highest weight $(6,0,-6)$ has
dimension three, not one.
\end{proposition}
\begin{proof}
The compact $U(3)$ half-sum of positive roots is $(1,0,-1)$.
Adding it to $(e-1,0,1-e)$ gives $(e,0,-e)$.
K17-U3-COVERAGE states this shift explicitly; K17-U3-HECKE
explicitly uses highest weight $(5,0,-5)$.
The dimension assertions follow independently from
KT-U3-DIM below. Thus the prompt's reference to weight $(6,0,-6)$
must be read as a parameter reference when discussing the original
$K_{12}$ certificate.
\end{proof}

\subsection{An exact dimension formula in every weight}

Put $c_j=[x^j](1+x+\cdots+x^5)^3$ and, for $d\geq0$, set
\begin{align*}
 I(d)&=\sum_{\substack{0\leq j\leq\min(15,d)\\j\equiv d\ (6)}}
      c_j\binom{(d-j)/6+2}{2}^{\!2},\\
 d_0+1&=6q+s,\qquad d_0=\lfloor d/2\rfloor,\quad0\leq s<6,\\
 T(d)&=6q^2+2qs+s,\qquad C(d)=\mathbf1_{3\mid d},\\
 A(d)&=\frac{I(d)+3T(d)+2C(d)}6,
 \qquad A(-1)=0.
\end{align*}

\begin{theorem}[KT-U3-DIM]
For every $d\geq0$,
\[
       \dim M_d=D(d):=A(d)-A(d-1).
\]
Equivalently, for $d=12n+r$, $n\geq0$, $0\leq r<12$, the
dimension is the following cubic quasipolynomial:
\[
\begin{array}{c|l@{\qquad}c|l}
r&D(12n+r)&r&D(12n+r)\\\hline
0&8n^3+2n^2+2n+1&6&(4n+3)(2n^2+2n+1)\\
1&2n(4n^2+2n+1)&7&(2n+1)(4n^2+6n+3)\\
2&8n^3+6n^2+4n+1&8&8n^3+18n^2+16n+5\\
3&(2n+1)(4n^2+2n+1)&9&2(n+1)(4n^2+6n+3)\\
4&(4n+1)(2n^2+2n+1)&10&8n^3+22n^2+22n+7\\
5&(2n+1)^3&11&8(n+1)^3
\end{array}
\]
In particular, $\dim M_5=1$, $\dim M_6=3$, and for every
positive $e\equiv0\pmod6$,
\[
                    \dim M_{e-1}=(e/6)^3.
\]
\end{theorem}
\begin{proof}
Use the harmonic realization of $V_d$ in bidegree $(d,d)$ in
three variables $z_i$ and their independent dual variables $w_i$.
The contraction map
$\Delta=\sum_i\partial_{z_i}\partial_{w_i}$ onto bidegree
$(d-1,d-1)$ is surjective, and its kernel is $V_d$.
This realization and its surjectivity in every degree were proved in
K17-HARMONIC, \texttt{harmonic\_hecke\_proof.md}, Section~1;
we use that result rather than reprove it.
Taking $\Gamma$-invariants is exact, since the characteristic is zero.
It therefore suffices to compute the invariant dimension in the full
bidegree space.

A monomial $z^a w^b$ is invariant under $\mu_6^3$ exactly when
$a_i\equiv b_i\pmod6$ for all $i$.
The remaining $S_3$ action permutes the three pairs $(a_i,b_i)$.
Burnside's lemma counts their orbits by averaging over the identity,
the three transpositions, and the two three-cycles.
For the identity, write $a_i=r_i+6u_i$ and $b_i=r_i+6v_i$ with
$0\leq r_i<6$. If $\sum r_i=j$, both $\sum u_i$ and $\sum v_i$
equal $(d-j)/6$, giving precisely $I(d)$ fixed monomials.
For a transposition the pairs in the two swapped positions agree.
Writing the exponents in these positions as $u,v$, the remaining
exponents are $d-2u,d-2v$. Thus $0\leq u,v\leq\lfloor d/2\rfloor$
and $u\equiv v\pmod6$; the congruence for the remaining position
then follows automatically. Among $d_0+1=6q+s$ integers, $s$
residue classes have size $q+1$ and $6-s$ have size $q$.
The count is $s(q+1)^2+(6-s)q^2=T(d)$.
A three-cycle fixes one monomial exactly when $3\mid d$, with
all six exponents equal to $d/3$. This gives $C(d)$.
The Burnside average is $A(d)$, and exactness gives $D(d)$.

For the table, substitute $d=12n+r$ in this proved finite formula.
One may sum over all $0\leq j\leq15$ with $j\equiv d\pmod6$:
the added terms have $(d-j)/6=-1$ or $-2$ and their polynomial
binomial coefficients vanish. Hence this substitution is a polynomial
identity, valid even at $n=0$, rather than an interpolation assumption.
Subtracting the adjacent $A$ polynomials gives the displayed table.
The last assertion is its $r=5,11$ cases.
\end{proof}

\paragraph{Scope.}
This is the total dimension. The endoscopic construction is supplied
separately in KT-U3-ENDO-INJECT; in particular KT-U3-ENDO-SIX
identifies the entire space at parameter $(6,0,-6)$ by its global
packet, in addition to the independent character computation below.

\subsection{A second method for the value at seven}

Write $\pi=3+\omega$ and $\bar\pi=2-\omega$.
For $r=1,2,3$, let $J_r$ be the $3$ by $3$ matrix whose first
$r$ by $r$ block consists of ones and whose other entries vanish.
Define
\[
 U_r=\pi I+\frac{\omega^{r-1}\bar\pi-\pi}{r}J_r,
 \qquad g_r=U_r/\pi.
\]
These are the three integral matrices
\[
\begin{aligned}
U_1&=\operatorname{diag}(2-\omega,3+\omega,3+\omega),\\[3pt]
U_2&=\begin{pmatrix}2+2\omega&-1+\omega&0\\
-1+\omega&2+2\omega&0\\0&0&3+\omega\end{pmatrix},
\\[3pt]
U_3&=\begin{pmatrix}1&-2-\omega&-2-\omega\\
-2-\omega&1&-2-\omega\\-2-\omega&-2-\omega&1\end{pmatrix}.
\end{aligned}
\]

\begin{theorem}[KT-U3-SECOND]
The scaled minuscule Hecke eigenvalue on $M_5$ is
\[
                         7^5T_7=157899.
\]
Equivalently, the normalized Satake trace is $157899/7^6$.
This is obtained by a character average on $\Gamma$, using the
three rank-one reflections above, independently of the harmonic
polynomial and neighbor data in the K17 computation.
\end{theorem}
\begin{proof}
The rank-one orthogonal projector onto the line spanned by
$(1,\ldots,1,0,\ldots,0)$ with $r$ ones is $J_r/r$.
Thus $U_r$ has eigenvalue $\omega^{r-1}\bar\pi$ on that line
and eigenvalue $\pi$ on its orthogonal complement.
Both eigenvalues have norm seven, so $U_r^*U_r=7I$.
The three displayed quotients are integral. Modulo $\pi$ they
have rank one, with image the indicated support line.
The elementary-divisor argument of K17-U3-HECKE, equations
(UH.7)--(UH.8), therefore identifies $g_r$ with the negative
minuscule correspondence at that line.

Reduction identifies $\mu_6$ with $\mathbf F_7^\times$.
Consequently $\Gamma$ acts on $\mathbf P^2(\mathbf F_7)$ with
exactly three orbits, determined by the number $r$ of nonzero
coordinates, of sizes
\[
             m_1=3,\qquad m_2=18,\qquad m_3=36.
\]
Every line has a representative with nonzero coordinates in
$\mu_6$, and conjugating $g_r$ by a corresponding monomial
unitary matrix gives its unitary reflection.
Distinct lines give distinct right cosets. Thus these independently
constructed matrices and their conjugates exhaust all $57$ cosets.

Let $P=|\Gamma|^{-1}\sum_{\gamma\in\Gamma}\rho_d(\gamma)$.
It is the projector onto $V_d^\Gamma$.
Because $P$ commutes with $\Gamma$, all conjugates of $g_r$ have
the same value of $\operatorname{tr}(P\rho_d(g_r))$.
The trace of the Hecke operator on $M_d$ is therefore
\[
 \operatorname{tr}(T_7\mid M_d)=
 \frac1{1296}\sum_{r=1}^3m_r
       \sum_{\gamma\in\Gamma}\chi_d(g_r\gamma).
\]
For a matrix $A$ let $h_j(A)$ denote the complete homogeneous
symmetric polynomial of degree $j$ in its eigenvalues, and put
$h_{-1}=0$.
The exact contraction sequence for the harmonic representation gives,
for a unitary $g$,
\[
 \chi_d(g)=h_d(g)\overline{h_d(g)}
                 -h_{d-1}(g)\overline{h_{d-1}(g)}.
\]
For $A=U_r\gamma$, multiplication by $7^d$ consequently gives
the integral expression
\[
       7^d\chi_d(g_r\gamma)
           =N(h_d(A))-7N(h_{d-1}(A)).                 \tag{KT-U.1}
\]
This computation needs no invariant harmonic polynomial and no
interpolation of its values.

Here is the complete compressed arithmetic certificate for $d=5$:
\[
\begin{array}{c|r|r|r}
r&m_r&\displaystyle\sum_{\gamma\in\Gamma}
       (N(h_5(U_r\gamma))-7N(h_4(U_r\gamma)))
       &\text{sum}/1296\\\hline
1&3&21781872&16807\\
2&18&5270832&4067\\
3&36&1233792&952
\end{array}
\]
Every entry is reproduced by the following fully specified recurrence.
For each of the six permutations $\sigma$ and each
$(u_1,u_2,u_3)\in\mu_6^3$, form $A_{ij}=(U_r)_{i,\sigma(j)}u_j$.
Set $s_1=\operatorname{tr}A$,
$s_2=\sum_{i<j}(A_{ii}A_{jj}-A_{ij}A_{ji})$, and $s_3=\det A$.
Starting at $h_0=1$ and $h_j=0$ for $j<0$, compute
\[
             h_j=s_1h_{j-1}-s_2h_{j-2}+s_3h_{j-3}
             \quad(1\leq j\leq5).
\]
Use $N(a+b\omega)=a^2-ab+b^2$ in (KT-U.1) and sum.
The standalone \texttt{character\_trace.py} implements exactly
these integer operations and no imported K17 computation or data.
This is a finite exact computation, not a Lean kernel certificate.
Since $\dim M_5=1$, its trace is its sole eigenvalue, namely
\[
           3(16807)+18(4067)+36(952)=157899.
\]
Complex conjugation exchanges the two minuscule correspondences
and preserves this one-dimensional coefficient space, so they have
the same eigenvalue. Gross, \emph{On the Satake isomorphism},
Section~3, equations (3.13)--(3.14), gives $T_7^+=7t_7$ under
the maximal-compact-volume-one convention. This proves the stated
Satake normalization.
\end{proof}

\subsection{An independent reflection calculation at thirteen}

\begin{theorem}[KT-U3-SECOND13]
The independent character-projector method at thirteen gives
\[
          13^3T_{13}|_{M_3}=54291,
          \qquad 13^5T_{13}|_{M_5}=6110859.
\]
The same calculation determines the trace of $13^dT_{13}$ in every
coefficient degree $d$ by the finite formula below.
\end{theorem}
\begin{proof}
Put $\pi=4+\omega$, $\bar\pi=3-\omega$, and $\delta=1-\omega$.
For each row of the following table set
\[
      U=\pi I+\frac{\eta\bar\pi-\pi}{n}vv^*,
      \qquad n=v^*v,\qquad g=U/\pi.
\]
\[
\begin{array}{c|c|r|r|r|r}
v&\eta&n&\text{orbit size}&b_3&b_5\\\hline
(1,0,0)&1&1&3&2197&371293\\
(1,1,0)&-\omega^2&2&18&1417&129233\\
(1,\delta,0)&\omega^2&4&18&-143&38233\\
(1,1,1)&\omega&3&36&532&7048\\
(1,1,\delta)&-\omega&5&108&52&16008
\end{array}
\]
The coefficient $(\eta\bar\pi-\pi)/n$ is respectively
$-1-2\omega$, $\omega$, $-2-\omega$, $-1+\omega$, and
$-1-\omega$, so all five numerator matrices are integral.
The rank-one projector $vv^*/n$ shows that the eigenvalues of
$U$ are $\eta\bar\pi,\pi,\pi$. Consequently
$U^*U=13I$ and $\det U=\eta\pi^2\bar\pi$.
Its reduction modulo $\pi$ has rank one and image the reduction
of $v$. The same local elementary-divisor argument as in
KT-U3-SECOND therefore places each $g$ in the required minuscule
double coset.

Here is the complete orbit classification. Modulo $\pi$ we have
$\omega=-4$ in $\mathbf F_{13}$, and the six Eisenstein units
are exactly the six nonzero squares. The residue of $\delta$ is
$5$, a nonsquare. Support-one lines form one orbit of size three.
For support two, the square class of the ratio of the two nonzero
coordinates gives two orbits, each of size $3\cdot6=18$;
interchanging the two coordinates preserves that square class.
For support three, consider the three square classes up to common
scaling. They are either all equal or have one different class.
The all-equal class gives an orbit of size $6^2=36$.
Permuting which coordinate has the different class gives one orbit
of size $3\cdot6^2=108$. Within each pattern the independent unit
multiplications are transitive. These five disjoint orbits have
total size $183=13^2+13+1$, so they exhaust the minuscule cosets.
The program also enumerates every projective line and checks this
partition exactly.

For each numerator $U$ put
\[
 b_d(U)=\frac1{1296}\sum_{\gamma\in\Gamma}
       \bigl(N(h_d(U\gamma))-13N(h_{d-1}(U\gamma))\bigr),
 \qquad h_{-1}=0.
\]
The character identity and scaling argument of KT-U3-SECOND
apply verbatim with $7$ replaced by $13$. Thus
$\operatorname{tr}(13^dT_{13}|M_d)$ is the sum of these five
$b_d$ values weighted by their orbit sizes. The $d=3,5$ values
from the complete-homogeneous recurrence are displayed in the
last two columns. Their weighted sums are respectively $54291$
and $6110859$. Since $D(3)=D(5)=1$, they are the eigenvalues.
The exact program \texttt{character\_trace13.py} imports only
the new character arithmetic from this audit and no archived
neighbor or operator data. It saves the five numerator matrices,
all 183 image lines, and all character sums through degree thirty.
\end{proof}

\subsection{A finite uniform formula for the entire spectrum}

\begin{proposition}[KT-U3-SPECTRUM]
The three reflections give a uniform exact formula for the
characteristic polynomial of $7^dT_7$ on $M_d$ for every $d\geq0$.
It does not give a uniform upper bound on the prime divisors of the
corresponding comparison integer.
\end{proposition}
\begin{proof}
For each projective line $L\in\mathbf P^2(\mathbf F_7)$ choose its
representative $v_L$ with first nonzero coordinate one and all
nonzero coordinates in $\mu_6$. If it has $r$ nonzero entries, put
\[
 U_L=\pi I+\frac{\omega^{r-1}\bar\pi-\pi}{r}v_Lv_L^*.
\]
The preceding proof gives a complete fixed list of $57$ integral
matrices, independent of $d$, each satisfying $U_L^*U_L=7I$.
The list is saved in \texttt{minuscule7\_cosets.json}; the independent
\texttt{generate\_cosets.py} checks every Gram identity, determinant,
and distinct projective image line.
Set $D=D(d)$. For $j\geq1$ and a sequence $L_1,\ldots,L_j$,
write $A=U_{L_1}\cdots U_{L_j}\gamma$.
The projector trace identity, applied to $T_7^j$, gives the exact
power sum
\[
 S_j=\operatorname{tr}((7^dT_7)^j\mid M_d)
 =\frac1{1296}
   \sum_{L_1,\ldots,L_j}\sum_{\gamma\in\Gamma}
       \bigl(N(h_d(A))-7^jN(h_{d-1}(A))\bigr).       \tag{KT-U.2}
\]
Indeed the product matrix satisfies $A^*A=7^jI$ and the same
homogeneity calculation as in (KT-U.1) applies.
Each $h_i$ is calculated by the three-term recurrence above.
Starting with $c_0=1$, define
\[
           c_j=-\frac1j\sum_{i=1}^j c_{j-i}S_i
           \quad(1\leq j\leq D).
\]
Newton's identities show that
$P_d(X)=X^D+c_1X^{D-1}+\cdots+c_D$ is the required
characteristic polynomial. Thus its roots, divided by $7^{d+1}$,
give all normalized Satake traces, with multiplicity, uniformly
in $d$. This is an exact finite formula, though its direct evaluation
is inefficient as $D$ grows.
The scaled action preserves the integral harmonic-polynomial
lattice over $O_E$, so $P_d\in O_E[X]$.
Formula (KT-U.2) puts every $S_j$ in $\mathbf Q$, and Newton's
identities therefore put $P_d$ in $\mathbf Q[X]$.
Since $O_E\cap\mathbf Q=\mathbf Z$, one has
$P_d\in\mathbf Z[X]$.
Complex conjugation exchanges the two minuscule correspondences
and conjugates their matrix coefficients. Their characteristic
polynomials are consequently equal. Thus the same polynomial
works for the positive operator required by the Satake formula.
No step in this construction bounds the prime factors of
$P_{e-1}(7^{2e}+7^e+1)$ independently of $e$.
\end{proof}

\paragraph{Check record and boundaries.}
The character program additionally verifies every reflection Gram
identity and rank-one reduction. Through degree $30$, its full
group-character averages agree with KT-U3-DIM, and independent
enumeration of exponent-pair orbits agrees with $A(d)$.
These finite comparisons check the implementation; the all-degree
dimension statement rests on the Burnside proof.
All-degree arithmetic prime elimination remains a separate question.
The identities $157899=7^6+40250$ and
$6110859=13^6+1284050$ have a global explanation in
KT-U3-ENDO-SIX: its proof first constructs the endoscopic packet
from the non-CM level-three weight-thirteen pair and then uses
the exact dimension to identify the compact form.
